\documentclass[../../main.tex]{subfiles}
\begin{document}
\section{带边子流形}\label{section:ch05-带边子流形}

设 $M$ 是(不带边的)光滑流形。一个 $M$ 的子集 $S\subseteq M$ 称为 $M$ 的\textbf{带边嵌入子流形},如果 $S$ 是带边拓扑流形并被赋予光滑结构,使包含映射 $S\hookrightarrow M$ 是光滑嵌入。注意,带边嵌入子流形的内部本身是 $M$ 的(不带边)嵌入子流形。余维数定义如常。

回忆边界坐标的定义,我们说子集 $S\subseteq M$ 满足\textbf{局部半空间 $k$-切片条件},如果 $S$ 的每一点都含于 $M$ 的某光滑坐标卡 $(U,\varphi)$ 的定义域内,使 $\varphi(S\cap U)$ 是 $\varphi(U)\cap\mathbb{H}^k$ 与某 $k$-切片
\begin{align*}
\bigl\{\bigl(x^1,\dots,x^k,x^{k+1},\dots,x^n\bigr):x^{k+1}=c^{k+1},\dots,x^n=c^n\bigr\}
\end{align*}
的交。任何这样的坐标卡称为 $S$ 的半空间切片坐标卡。

下面的命题总结了带边嵌入子流形最常用的例子。

\begin{proposition}[Level Sets with Boundary]\label{proposition:ch05-level-sets-boundary}
设 $M$ 是(不带边的)光滑流形, $N$ 是(不带边的)光滑流形, $\Phi:M\to N$ 是光滑映射, $S\subseteq M$ 是正则水平集 $\Phi^{-1}(c)$。若 $D\subseteq M$ 是带边嵌入子流形,且 $\partial D$ 与 $S$ 横截相交(见\cref{section:ch06-横截性}),则 $S\cap D$ 是 $D$ 的带边嵌入子流形,其边界为 $S\cap\partial D$。
\end{proposition}

\begin{proof}
见习题 5-13。
\end{proof}

下面是带边嵌入子流形的局部切片刻画,它与\cref{theorem:ch05-slice-criterion} 平行。

\begin{theorem}[Local Slice Criterion for Submanifolds with Boundary]\label{theorem:ch05-slice-boundary}
设 $M$ 是(不带边的)光滑流形。若 $S\subseteq M$ 是带边嵌入 $k$ 子流形,则 $S$ 满足局部半空间 $k$-切片条件。反之,若 $S\subseteq M$ 是满足局部半空间 $k$-切片条件的子集,则在子空间拓扑下, $S$ 是带边拓扑 $k$ 流形,并具有一个光滑结构使它成为 $M$ 的真嵌入子流形。
\end{theorem}

\begin{proof}
完全仿照\cref{theorem:ch05-slice-criterion},只是在构造 $S$ 的坐标卡时使用 $\mathbb{H}^k$ 的坐标,并验证转移映射是 $\mathbb{H}^k$ 开子集之间的光滑映射。细节留作习题 5-14。
\end{proof}

\begin{example}\label{example:ch05-hemisphere-submanifold}
令
\begin{align*}
D=\bigl\{(x,y,z)\in\mathbb{R}^3:x^2+y^2+z^2=1,\ z\geqslant 0\bigr\}
\end{align*}
为单位球面的上半部分。它是 $S^2$ 与闭半空间 $z\geqslant 0$ 的交,故它是带边嵌入子流形(见习题 5-15)。其边界为赤道圆 $\{(x,y,0):x^2+y^2=1\}$。
\end{example}

若 $M$ 本身就是带边流形,则还有另外几种方式使子集 $S\subseteq M$ 值得研究。例如,若 $M$ 是带边光滑 $n$ 流形,则它的边界 $\partial M$ 本身是 $M$ 的一个 $(n-1)$ 维嵌入子流形(不带边)。更一般地,若 $M$ 是带边光滑流形, $S\subseteq M$ 是子集,我们称 $S$ 是 $M$ 的\textbf{(带边或不带边)子流形},如果 $S$ 在由 $M$ 诱导的拓扑下是带边或不带边流形,并被赋予光滑结构使包含映射是光滑嵌入。然而,由于 $M$ 的边界, $S$ 可能以各种方式与 $\partial M$ 相交,因此我们不在此一般情形下证明任何结果。

\begin{figure}[ht]
\centering
\includegraphics[width=0.8\textwidth]{fig-05-14.png}
\caption{带边流形中的带边子流形}\label{figure:ch05-submanifold-boundary}
\end{figure}

\end{document}